\subsection{Problem Setup}

We formalize the problem of kinodynamic motion planning using edge bundles similarly to \cite{shome2021asymptotically}.
The robot state space $\mathcal{Q} \subset \mathbb{R}^m$ is defined where $q \in \mathcal{Q}$ represents the joint angles of an $m$-dimensional robot. The control space $\mathcal{U} \subset \mathbb{R}^m$ contains elements $\dot{q} \in \mathcal{U}$ that represent joint velocities.
An \textsl{edge bundle} is a set of disconnected edges generated by applying randomly sampled control parameters to random robot states for random time durations, thereby connecting random start states to resultant end states:
\begin{center}
$\mathcal{E} = \{(q_0, \boldsymbol{\dot{q}}, \Delta t, q_f) \mid q_0, q_f \in \mathcal{Q}, \hspace{3pt} \boldsymbol{\dot{q}} \in \mathcal{U}\}$
\end{center}
such that $f(q_0, \boldsymbol{\dot{q}}) = q_f$, where $f(\cdot)$ defines the forward dynamics model of the robot and $\boldsymbol{\dot{q}} = (\dot{q}_0, \cdots, \dot{q}_{f-1})$ represents a sequence of joint velocities. The $i$-th edge is denoted as $e_i$ and defined as the tuple $(q_0^i, \boldsymbol{\dot{q}^i}, \Delta t^i, q_f^i)$.
The objective of \textsl{LazyBoE} is to find a continuous, collision-free trajectory
$$\pi(t) : [0,1] \rightarrow \mathcal{Q}_{\text{free}} \subseteq \mathcal{Q}, \quad \pi(0) = q_0, \quad \pi(1) = q_f$$
through the state and control space discretization enabled by $\mathcal{E}$, while respecting the dynamics constraints $\mathcal{D}$ of the robot.
